\documentclass[11pt]{article}
\usepackage[margin=1in]{geometry}
\usepackage{amsmath,amssymb,amsthm}
\usepackage{graphicx}
\usepackage[hidelinks]{hyperref}

\newtheorem{theorem}{Theorem}
\newtheorem{lemma}{Lemma}
\newtheorem{definition}{Definition}
\newtheorem{remark}{Remark}

\newcommand{\N}{\mathbb N}
\newcommand{\R}{\mathbb R}
\newcommand{\one}{\mathbf 1}
\newcommand{\omegaone}{\omega_1}
\newcommand{\clspan}{\overline{\operatorname{span}}}
\newcommand{\supp}{\operatorname{supp}}

\title{A countable-block overcomplete set for an ordinal interval space}
\author{Agent lane 16}
\date{2026-07-03}

\begin{document}
\maketitle

\noindent\textbf{Status.}
Candidate partial result, likely valid.  This packet proves the first
countable-block case left open by the earlier lane-16 packet: in ZFC,
\(C([0,\omega_1\cdot\omega])\) admits an overcomplete set.  It does not claim
the full \(\xi<\omega_2\) solution to Koszmider's Question 48.

\section{Source question}

The source paper is:

\begin{quote}
P. Koszmider, \emph{On the existence of overcomplete sets in some classical
nonseparable Banach spaces}, arXiv:2006.00806.
\end{quote}

On page 25, Question 48 asks:

\begin{quote}
Can one prove in ZFC that the Banach spaces \(C([0,\xi])\) for all ordinals
\(\xi<\omega_2\) admit overcomplete sets?
\end{quote}

\begin{center}
\includegraphics[width=0.92\linewidth]{figures/open_problem_crop.png}
\end{center}

\section{Definitions and source tools}

\begin{definition}
Let \(X\) be a Banach space. A set \(D\subseteq X\) is \emph{overcomplete} if
\(|D|=\operatorname{dens}(X)\) and every subset \(A\subseteq D\) with
\(|A|=|D|\) has dense linear span in \(X\).
\end{definition}

We use two standard facts used explicitly in Koszmider's paper.  First,
because \([0,\omega_1]\) is scattered, every functional on
\(C([0,\omega_1])\) is an atomic measure with countable support.  Second,
there are almost coherent injections \(e_\alpha:\alpha\to\N\)
\((\alpha<\omega_1)\) such that
\[
   e_\beta\!\upharpoonright\alpha =^* e_\alpha
   \qquad(\alpha<\beta<\omega_1),
\]
and they may be chosen with \(\N\setminus e_\alpha[\alpha]\) infinite.
As in the source proof, this implies the following exact stabilization:
if \(A\subseteq\omega_1\) is uncountable and \(\gamma<\omega_1\), then there
are an uncountable \(A'\subseteq A\setminus(\gamma+1)\) and an injection
\(g:\gamma+1\to\N\) such that
\[
   e_\alpha\!\upharpoonright(\gamma+1)=g
   \qquad(\alpha\in A').
\]

\section{Proof Intuition}

The previous partial packet stopped before \(\omega_1\cdot\omega\) because the
finite-codimensional countable-support criterion does not survive countably
many \(\omega_1\)-blocks.  The obstacle is cancellation: a functional on a
countable \(c_0\)-sum is a summable sequence of block functionals, and the
block evaluations can cancel each other.

The construction below prevents that cancellation by assigning each block its
own sparse set of Taylor exponents.  On an uncountable subfamily, the coherent
injections make the value of each fixed block functional an analytic power
series in one real parameter \(r_\alpha\).  Since different blocks use
disjoint exponent sets, an analytic identity forces the coefficients of every
block separately to vanish.  A small endpoint-detector series with unbounded
coefficients then kills the endpoint mass in each copy of \(C([0,\omega_1])\),
just as in the strengthened \(C([0,\omega_1])\) construction from the earlier
run work.

\section{The countable-block theorem}

\begin{theorem}\label{thm:model}
Let
\[
   Y=C([0,\omega_1]).
\]
Then the Banach space
\[
   X=\R\oplus \Bigl(\bigoplus_{n<\omega}Y\Bigr)_{c_0}
\]
admits a set \(\{x_\alpha:\alpha<\omega_1\}\) such that every nonzero
functional on \(X\) vanishes on at most countably many \(x_\alpha\)'s.  In
particular, \(X\) admits an overcomplete set.
\end{theorem}

\begin{proof}
Write \(K=[0,\omega_1]\), \(u_\alpha=\one_{[0,\alpha]}\in C(K)\)
\((\alpha<\omega_1)\), and \(\one=\one_K\).  In this proof
\(\N=\{0,1,2,\ldots\}\).

Choose pairwise disjoint infinite sets
\[
   P_\infty,P_0,P_1,\ldots\subseteq\N
\]
and enumerations \(p_j:\N\to P_j\), \(j\in\{\infty,0,1,\ldots\}\), such that
\[
   p_j(k)\geq 2k+1\qquad(k<\omega).
\]
Put
\[
   q_j(t)=\sum_{k=0}^{\infty}2^k t^{p_j(k)}
   \qquad (|t|<1/\sqrt2).
\]
Fix \(0<b<1/2\), and choose pairwise distinct \(r_\alpha\in(b/2,b)\) for
\(\alpha<\omega_1\).  Also fix positive scalars \(a_n\downarrow0\), for
instance \(a_n=2^{-n-1}\).

For \(j<\omega\) and \(\alpha<\omega_1\), define \(v_\alpha^{(j)}\in Y\) by
\[
 v_\alpha^{(j)}
   = r_\alpha^{p_j(0)}u_\alpha
     +\sum_{\beta<\alpha}
        r_\alpha^{p_j(e_\alpha(\beta)+1)}(u_\beta-u_\alpha)
     +q_j(r_\alpha)\one .
\]
The series converges absolutely in \(Y\), because \(e_\alpha\) is injective and
\(\sum_k b^{p_j(k)}\leq\sum_k b^{2k+1}<\infty\).  The vectors
\(v_\alpha^{(j)}\) are uniformly bounded in \(j\) and \(\alpha\).  Hence
\[
   x_\alpha
      =\bigl(q_\infty(r_\alpha), (a_n v_\alpha^{(n)})_{n<\omega}\bigr)
      \in X
\]
is well-defined.  Since \(q_\infty\) has positive coefficients, it is strictly
increasing on \((0,b)\), so the \(x_\alpha\)'s are pairwise distinct.

We prove the functional separation property.  Let
\[
   \Phi=(c,(\mu_n)_{n<\omega})\in
      \R\oplus\Bigl(\bigoplus_{n<\omega}Y^*\Bigr)_{\ell_1}=X^*
\]
and suppose that \(\Phi(x_\alpha)=0\) for all \(\alpha\) in some uncountable
set \(A\subseteq\omega_1\).  For each \(n\), the measure \(\mu_n\) has
countable support in \(K\); since there are only countably many \(n\)'s, the
union of the supports below \(\omega_1\) is bounded.  Choose
\(\gamma<\omega_1\) such that
\[
   \bigcup_{n<\omega}\bigl(\supp(\mu_n)\setminus\{\omega_1\}\bigr)
       \subseteq [0,\gamma].
\]
By exact stabilization, pass to an uncountable
\(A'\subseteq A\setminus(\gamma+1)\) and an injection \(g:\gamma+1\to\N\) such
that \(e_\alpha\!\upharpoonright(\gamma+1)=g\) for every \(\alpha\in A'\).

For \(n<\omega\), put
\[
   b_{n,\beta}=\mu_n(u_\beta)\quad(\beta\leq\gamma),\qquad
   d_n=\mu_n(\one),
\]
and let \(s_n\) be the eventual value of \(\mu_n(u_\xi)\) for
\(\xi>\gamma\).  For \(\alpha\in A'\), with \(t=r_\alpha\), the definition of
\(v_\alpha^{(n)}\) gives
\[
 \mu_n(v_\alpha^{(n)})
   =s_n t^{p_n(0)}
    +\sum_{\beta\leq\gamma}(b_{n,\beta}-s_n)
       t^{p_n(g(\beta)+1)}
    +d_n q_n(t).
\]
Consequently the analytic function
\[
   F(t)=c\,q_\infty(t)+
        \sum_{n=0}^{\infty} a_n\left[
        s_n t^{p_n(0)}
        +\sum_{\beta\leq\gamma}(b_{n,\beta}-s_n)
           t^{p_n(g(\beta)+1)}
        +d_n q_n(t)\right]
\]
vanishes at every \(r_\alpha\) with \(\alpha\in A'\).  The series defining
\(F\) converges normally on compact subdiscs of \(|t|<1/\sqrt2\), because
\((\|\mu_n\|)_n\in\ell_1\), the \(a_n\)'s are bounded, and the displayed
series for each block is bounded by a constant depending only on \(b\) times
\(\|\mu_n\|\).  The set \(\{r_\alpha:\alpha\in A'\}\) has an accumulation
point in \((b/2,b)\), so the identity theorem gives \(F\equiv0\).

Compare Taylor coefficients.  The sets \(P_\infty,P_0,P_1,\ldots\) are
pairwise disjoint.  The coefficients on \(P_\infty\) are \(c\,2^k\); hence
\(c=0\).  Fix \(n<\omega\).  Since \(a_n>0\), the coefficients on \(P_n\)
vanish.  The range of \(g\) has infinite complement, so there are infinitely
many \(k\)'s outside
\[
   \{0\}\cup\{g(\beta)+1:\beta\leq\gamma\}.
\]
For such \(k\), the coefficient of \(t^{p_n(k)}\) is \(a_n d_n2^k\), so
\(d_n=0\).  The coefficient of \(t^{p_n(0)}\) then gives \(s_n=0\).  Finally,
for each \(\beta\leq\gamma\), the coefficient of
\(t^{p_n(g(\beta)+1)}\) gives \(b_{n,\beta}=0\).

Thus \(\mu_n(\one)=0\), \(\mu_n(u_\beta)=0\) for all \(\beta\leq\gamma\), and
\(\mu_n(u_\xi)=s_n=0\) for all \(\xi>\gamma\).  The linear span of
\(\one\) and the \(u_\xi\)'s is dense in \(C([0,\omega_1])\): it is a unital
algebra separating points of the compact ordinal interval.  Therefore
\(\mu_n=0\).  Since this holds for every \(n\), and \(c=0\), we have
\(\Phi=0\).

We have proved that any functional vanishing on an uncountable subfamily of
\(\{x_\alpha:\alpha<\omega_1\}\) must be zero.  Since \(\operatorname{dens}X
=\omega_1\), it follows that every uncountable subfamily has dense linear
span by Hahn--Banach.  Hence the family is overcomplete in \(X\).
\end{proof}

\begin{theorem}\label{thm:ordinal}
In ZFC, \(C([0,\omega_1\cdot\omega])\) admits an overcomplete set.
\end{theorem}

\begin{proof}
Let \(K=[0,\omega_1\cdot\omega]\), and put \(t_\infty=\omega_1\cdot\omega\).
For \(n<\omega\), let
\[
   I_0=[0,\omega_1],
   \qquad
   I_n=(\omega_1\cdot n,\omega_1\cdot(n+1)]\quad(n\geq1).
\]
Then \(K\setminus\{t_\infty\}\) is the topological disjoint union of the
compact-open pieces \(I_n\), and each \(I_n\) is homeomorphic to
\([0,\omega_1]\).  The map
\[
   Tf=\Bigl(f(t_\infty),
       \bigl(f|_{I_n}-f(t_\infty)\one_{I_n}\bigr)_{n<\omega}\Bigr)
\]
is a linear isomorphism from \(C(K)\) onto
\[
   \R\oplus\Bigl(\bigoplus_{n<\omega} C([0,\omega_1])\Bigr)_{c_0}.
\]
Indeed, continuity at \(t_\infty\) is exactly the condition that the displayed
block sequence tends to zero in norm.  Apply Theorem~\ref{thm:model} and
pull the overcomplete set back through \(T\).
\end{proof}

\section{Limitations and novelty check}

This improves the earlier lane-16 packet for ordinal intervals below
\(\omega_1\cdot\omega\) by adding the first missing countable-block endpoint
\(\xi=\omega_1\cdot\omega\).

The proof is not yet a full answer to Question 48.  It handles a countable
\(c_0\)-sum of copies of \(C([0,\omega_1])\) by assigning countably many
blocks disjoint Taylor supports.  The same simple coefficient separation does
not directly handle ordinal intervals whose natural block structure involves
\(\omega_1\)-many interacting \(\omega_1\)-blocks, such as the
\(\omega_1^2\)-level.

Local indexes were searched for arXiv:2006.00806, ordinal-interval
overcomplete keywords, and the previous packet slug.  A fresh web/literature
pass checked the exact Question 48 wording, ordinal \(C([0,\alpha])\)
classification phrases, Bessaga--Pelczynski/Kislyakov/Gul'ko--Os'kin
classification references, and \(C([0,\omega_1\cdot\omega])\) isomorphism
phrases.  This found classification literature and related operator theory on
ordinal \(C(K)\) spaces, but no later paper explicitly resolving Question 48
or this countable-block subcase.

\section{Human review recommendation}

Review as a likely valid partial result.  The main points to check are:
\begin{enumerate}
\item the coefficient-separation argument using pairwise disjoint Taylor
      exponent sets \(P_n\);
\item the normal convergence of the analytic function \(F\);
\item the use of exact stabilization for one countable bound \(\gamma\) that
      simultaneously covers all block measures;
\item the identification of \(C([0,\omega_1\cdot\omega])\) with the scalar
      plus countable \(c_0\)-sum model.
\end{enumerate}

\begin{thebibliography}{9}

\bibitem{Koszmider}
P. Koszmider,
\emph{On the existence of overcomplete sets in some classical nonseparable
Banach spaces},
J. Funct. Anal. \textbf{281} (2021), no. 9, 109172;
\href{https://arxiv.org/abs/2006.00806}{arXiv:2006.00806}.

\bibitem{RussoSomaglia}
T. Russo and J. Somaglia,
\emph{Overcomplete sets in non-separable Banach spaces},
Proc. Amer. Math. Soc. \textbf{149} (2021), 701--714;
\href{https://arxiv.org/abs/1912.08690}{arXiv:1912.08690}.

\end{thebibliography}

\end{document}
